\hypertarget{block-sieve}{%
\section{The quantitative finite block sieve}\label{block-sieve}}

We prove the finite estimate used for the rectangle upper bound and the
one-dimensional sieves. Its constants are independent of the block
count.

\Needspace{6\baselineskip}

\hypertarget{block-sieve-statement}{%
\subsection{Statement}\label{block-sieve-statement}}

Let \(z>1\), let \(h\ge2\) be even, and let \(P_j\), \(0\le j<J\), be
finite sets of primes in \((z^{2^{-(j+1)}},z^{2^{-j}}]\). These blocks
are disjoint; put \(P=\bigcup_jP_j\). Let \(\mathcal X\) be a finite set
of objects with weights \(w(x)\ge0\) and bad-prime sets \(B(x)\).
Suppose \(X\ge0\), \(0\le g(p)\le1-\eta\) for \(p\in P\), where
\(\eta>0\), \(C\ge0\), and \(\sum_{p\in P_j}g(p)\le C\). For every
\(s\subseteq P\), assume

\[
\sum_{x\in\mathcal X}w(x)\prod_{p\in s}1_{p\in B(x)}
 =X\prod_{p\in s}g(p)+r\!\left(\prod_{p\in s}p\right). \tag{BS1}
\]

Write

\[
\begin{gathered}
V=\prod_{p\in P}(1-g(p)),\qquad D=z^{4h+2},\qquad
\mathcal E=\sum_{\substack{d\mid\prod_{p\in P}p\\d\le D}}|r(d)|,\\
K(C,\eta)=\frac{\exp(eC+C/\eta)}{1-e^{-1}},\qquad
\rho=K(C,\eta)e^{-h},\qquad
S=\sum_{x\in\mathcal X}w(x)1_{B(x)\cap P=\varnothing}.
\end{gathered}
\]

Then

\[
XV(1-\rho e^\rho)-\mathcal E
 \le S\le XV e^\rho+\mathcal E. \tag{BS2}
\]

Every divisor in \(\mathcal E\) is squarefree. Both \(J\) and \(h\) may
vary.

\Needspace{6\baselineskip}

\hypertarget{block-sieve-proof}{%
\subsection{Proof}\label{block-sieve-proof}}

Set \(k_j=(j+1)h\). For variables \(t_p\in[0,1]\), define

\[
\begin{aligned}
a_j(t)&=\prod_{p\in P_j}(1-t_p),\\
u_j(t)&=\sum_{n=0}^{k_j}(-1)^n
             \sum_{\substack{s\subseteq P_j\\|s|=n}}\prod_{p\in s}t_p,\\
e_j(t)&=\sum_{\substack{s\subseteq P_j\\|s|=k_j+1}}\prod_{p\in s}t_p.
\end{aligned}
\]

Even Bonferroni truncation gives \(a_j\le u_j\) and
\(0\le u_j-a_j\le e_j\). To verify it, take independent Bernoulli
variables of means \(t_p\) and average
\(\sum_{n=0}^k(-1)^n\binom{q}{n}=(-1)^k\binom{q-1}{k}\) for \(q>0\),
handling \(q=0\) separately. A telescoping expansion now gives

\[
\prod_j u_j-\sum_j e_j\prod_{i\ne j}u_i
 \le\prod_j a_j\le\prod_j u_j. \tag{BS3}
\]

Apply (BS3) with \(t_p=1_{p\in B(x)}\), multiply by \(w(x)\) and sum.
Both bounding polynomials have monomial coefficients of absolute value
at most one. A subset contributes to the upper polynomial only when its
degree in every block is at most \(k_j\). It contributes to the
subtracted sum only when exactly one block has degree \(k_j+1\) and all
other degrees are at most their cutoffs. These cases are disjoint, and
the exceptional block is unique.

Every surviving monomial has prime product at most \(D\): its exponent
of \(z\) is at most \(h\sum_{j\ge0}(j+1)2^{-j}+1=4h+1<4h+2\). Distinct
prime subsets have distinct products. Inserting (BS1) therefore costs at
most \(\mathcal E\) for each bounding polynomial, with no extra divisor
coefficients or factor \(J\).

Evaluate next at \(t_p=g(p)\). Abbreviate \(a_j=a_j(g)>0\). Since
\(\log(1-t)\ge-t/\eta\) for \(0\le t\le1-\eta\), \(a_j\ge e^{-C/\eta}\).
Moreover

\[
e_j(g)\le\frac{C^{k_j+1}}{(k_j+1)!}\le\exp(eC-k_j-1),
\]

by retaining one term in the series for \(\exp(eC)\). Thus
\(\delta_j=e_j(g)/a_j\) satisfies \(\sum_j\delta_j\le\rho\). From
\(a_j\le u_j(g)\le a_j(1+\delta_j)\) we obtain

\[
V\le\prod_j u_j(g)\le Ve^\rho,
\qquad
\sum_j e_j(g)\prod_{i\ne j}u_i(g)\le V\rho e^\rho.
\]

Together with (BS3) and the remainder bound, these inequalities prove
(BS2).

\Needspace{6\baselineskip}

\hypertarget{block-sieve-constants-in-the-applications}{%
\subsection{Constants in the
applications}\label{block-sieve-constants-in-the-applications}}

For \protect\hyperlink{rectangles}{rectangles}, the local estimates give
\(C=12\), \(\eta=2/9\), and \(h=2\). Hence
\(\rho=e^{12e+52}/(1-e^{-1})\) and \(D=z^{10}\), exactly as in (R7). The
block count may grow.

For the preliminary sieve on \(u+2\), the density is zero at \(2\) and
\(1/(p-1)\) at odd primes. The reciprocal-prime block bounds give
\(C=8\), and \(\eta=1/2\). Take \(h=2\lfloor1/(40b)\rfloor\). If
\(0<b\le1/1000\), then \(h\ge1/(40b)\), \(h\ge50\) and \(\rho\le1\).
Consequently

\[
\rho e^\rho\le C_s e^{-1/(40b)},\qquad C_s=3K(8,1/2).
\]

This is the constant in (S7). The level exponent satisfies
\(b(4h+2)\le1/5+2b\). For the conditioned sieve in (S8b), use the same
\(C,\eta\) and let an even depth tend to infinity while keeping
\(W^{4h+2}\) below the remaining level after conditioning on the prime.
Estimate (BS2) is uniform in this depth and the block count.

For ordinary rough integers the density is \(1/p\), with block bound
\(C=4\) and gap \(\eta=1/2\). Put \(K_0=K(4,1/2)\) and
\(C_{\rm sieve}=K_0e^{K_0}\). Since \(\rho=K_0e^{-h}\le K_0\), both
relative errors in (BS2) are at most \(C_{\rm sieve}e^{-h}\). Counting
multiples of each squarefree divisor has error at most one, so the total
error is at most \(2D\), as used in (S8d). Both \(C_s\) and
\(C_{\rm sieve}\) are fixed before \(b\) is chosen.
